\centerline{\bf \Huge Вокруг ортоцентра}

\problem{1}
Дан остроугольный треугольник $A B C ; A A^{\prime}, B B^{\prime}, C C^{\prime}$ - его высоты, $O$ - центр описанной окружности. Докажите, что касательная в точке $B$ к описанной окружности параллельна прямой $A^{\prime} C^{\prime}$, а $O B \perp A^{\prime} C^{\prime}$.

\problem{2}
Высоты треугольника $A B C$ пересекаются в точке $H$. Докажите, что радиусы описанных окружностей треугольников $A B C, A B H, B C H$ и $C A H$ равны.

\problem{3}
Вокруг остроугольного треугольника $A B C$ описана окружность. Продолжения высот треугольника, проведённых из вершин $A$ и $C$, пересекают окружность в точках $E$ и $F$ соответственно, $D$ - произвольная точка на меньшей дуге $A C, K$ - точка пересечения $D F$ и $A B, L$ - точка пересечения $D E$ и $B C$. Докажите, что прямая $K L$ проходит через ортоцентр треугольника $A B C$.

\problem{4}
Пусть $B H_b, C H_c$ - высоты треугольника $A B C$. Прямая $H_b H_c$ пересекает описанную окружность треугольника $A B C$ в точках $X$ и $Y$. Точки $P$ и $Q$ симметричны $X$ и $Y$ относительно $A B$ и $A C$ соответственно. Докажите, что $P Q \| B C$.

\problem{5}
Пусть $A_1, B_1, C_1$ - основания высот остроугольного треугольника $A B C$. Окружность, вписанная в треугольник $A_1 B_1 C_1$, касается сторон $A_1 B_1$, $A_1 C_1, B_1 C_1$ в точках $C_2, B_2, A_2$. Докажите, что прямые $A A_2, B B_2, C C_2$ пересекаются в одной точке, лежащей на прямой Эйлера треугольника $A B C$.


\problem{6}
Высоты $A A^{\prime}$ и $C C^{\prime}$ остроугольного треугольника $A B C$ пересекаются в точке $H$. Точка $B_0$ - середина стороны $A C$. Докажите, что точка пересечения прямых, симметричных $B B_0$ и $H B_0$ относительно биссектрис углов $A B C$ и $A H C$ соответственно, лежит на прямой $A^{\prime} C^{\prime}$.

\problem{7}
Пусть $H$ и $O$ - ортоцентр и центр описанной окружности треугольника $A B C$. Описанная окружность треугольника $A O H$, пересекает серединный перпендикуляр к $B C$ в точке $A_1$. Аналогично определяются точки $B_1$ и $C_1$. Докажите, что прямые $A A_1, B B_1$ и $C C_1$ пересекаются в одной точке.

\problem{8}
Окружность $\Omega$ с центром в точке $O$ описана около остроугольного треугольника $A B C$, в котором $A B<B C$; его высоты пересекаются в точке $H$. На продолжении отрезка $B O$ за точку $O$ отмечена точка $D$ такая, что $\angle A D C=\angle A B C$. Прямая, проходящая через точку $H$ параллельно прямой $B O$, пересекает меньшую дугу $A C$ окружности $\Omega$ в точке $E$. Докажите, что $B H=D E$.

